\subsection{结构的种类}

让 \(\mathfrak{G}\) 是一个比集合论更强的理论。在 \(\mathfrak{G}\) 中的结构种类是一个文本 \(\Sigma\) 由以下汇编组成：

\begin{enumerate}
\def\labelenumi{(\arabic{enumi})}
\item
  一定数量的字母 \(x_1, \ldots, x_n, s\) ，彼此不同，且与常数 \(\mathfrak{G}\) ; \(x_1, \ldots, x_n\) 被称为该结构物种的主基集 \(\Sigma\) ;
\item
  一定数量的项 \(A_{1}, \ldots, A_{m}\) 在 G 中，其中没有一个字母 \(x_{1}, \ldots, x_{n}, s\) 出现，并且被称为辅助基集 \(\Sigma\) ; \(\Sigma\) 可能不包含任何辅助基集（但必须至少包含一个主基集）；
\item
  一种典型化 \(\mathrm{T}\{x_1, \ldots, x_n, s\}\) :
\end{enumerate}

\[
s \in \mathrm{S} (x _ {1}, \dots , x _ {n}, \mathrm{A} _ {1}, \dots , \mathrm{A} _ {m}),
\]

其中 S 是定义在 \(n + m\) 术语（编号1）；T \(\{x_{1}, \ldots, x_{n}, s\}\) 被称为结构种类的典型刻画 \(\Sigma\) ;

\begin{enumerate}
\def\labelenumi{(\arabic{enumi})}
\setcounter{enumi}{3}
\tightlist
\item
  一个关系 \(R\{x_{1},\ldots,x_{n},s\}\) 该关系在类型化T下（在C中）是可迁移的， \(x_{i}\) 作为主基集和 \(A_{h}\) 辅助基集（第3号）；R被称为结构种类的公理 \(\Sigma\) .
\end{enumerate}

理论 \(C_{\Sigma}\) 具有与C相同的公理模式，其显式公理为C的公理，再加上公理“T和R”，称为结构物种的理论 \(\Sigma\) 。 \(C_{\Sigma}\) 的常量因此是C的常量以及出现在T或R中的字母。

让 \(\mathfrak{C}'\) 是比 \(\mathfrak{C}\) 更强的理论，并且让 \(\mathbf{E}_1, \ldots, \mathbf{E}_n\) ，U是 \(\mathfrak{C}'\) 中的项。在理论 \(\mathfrak{C}'\) 中，U被称为物种 \(\Sigma\) 在主基集上 \(\mathbf{E}_1, \ldots, \mathbf{E}_n\) ，以 \(\mathbf{A}_1, \ldots, \mathbf{A}_m\) 作为辅助基集，如果该关系

\[
" \mathrm{T} \left\{\mathrm{E} _ {1}, \dots , \mathrm{E} _ {n}, \mathrm{U} \right\} \text { 且 } \mathrm{R} \left\{\mathrm{E} _ {1}, \dots , \mathrm{E} _ {n}, \mathrm{U} \right\}"
\]

在 \(\mathfrak{C}'\) 中是定理。当这种情况发生时，对于每个定理 \(\mathbf{B}\{x_1, \ldots, x_n, s\}\) （在理论 \(\mathfrak{C}_{\Sigma}\) 中），关系 \(\mathbf{B}\{\mathrm{E}_1, \ldots, \mathrm{E}_n, \mathrm{U}\}\) 在 \(\mathfrak{C}'\) 中是定理（第一章，§2，第3号）。在 \(\mathfrak{C}_{\Sigma}\) 中，常量 \(s\) 被称为物种 \(\Sigma\) .

的泛型结构。在理论 \(\mathfrak{G}'\) 中，主基集 \(\mathbf{E}_1, \ldots, \mathbf{E}_n\) 被称为具有结构U。显然，U是集合

\[
\mathrm{S} (\mathbf {E} _ {1}, \dots , \mathbf {E} _ {n}, \mathbf {A} _ {1}, \dots , \mathbf {A} _ {m}).
\]

因此， \(\mathrm{S}(\mathbf{E}_{1},\ldots,\mathbf{E}_{n},\mathbf{A}_{1},\ldots,\mathbf{A}_{m})\) 中满足关系 \(R\{\mathbf{E}_{1},\ldots,\mathbf{E}_{n},\mathbf{V}\}\) 的元素 V 所成的集合，就是种类 \(\Sigma\) 在基集 \(E_{1},\ldots,E_{n}\) 上的所有结构之集（它可能是空集）。

示例

\begin{enumerate}
\def\labelenumi{(\arabic{enumi})}
\tightlist
\item
  取 \(\mathfrak{C}\) 为集合论，并考虑这样一种结构：它没有辅助基集，只有一个主基集 A，典型特征为 \(s \in \mathfrak{P}(A \times A)\) ，公理为
\end{enumerate}

\[
s \circ s = s \quad \text { 且 } \quad s \cap \overline{s} = \Delta_{\mathbf{A}}
\]

\((\Delta_{\mathbf{A}}\) 为对角线 \(\mathbf{A} \times \mathbf{A})\) ，这是关于类型化的一种可迁移关系 \(s \in \mathfrak{P}(\mathbf{A} \times \mathbf{A})\) ，这一点很容易验证。显然，这类结构的理论正是有序集的理论（第三章，§1，第3节）；因此，这样定义的结构类型也被称为 \(\mathbf{A}\) 上的序结构物种。在第三章中，我们看到了许多带有这种物种结构的集合的例子。

\begin{enumerate}
\def\labelenumi{(\arabic{enumi})}
\setcounter{enumi}{1}
\item
  取 \(\mathfrak{G}\) 为集合论，并考虑这样一种结构：它没有辅助基集，只有一个主基集 A，典型特征为 \(F \in \mathfrak{P}((A \times A) \times A)\) ，公理是可迁移关系“F 是一个定义域为 \(A \times A\) 的函数图”。这种结构是所谓代数结构的特例；以 F 为图的函数（从 \(A \times A\) 到 A 的映射），称为这种结构的（处处有定义的）内部合成律。
\item
  和之前一样，令 \(\mathfrak{C}\) 设集合论为理论，并考虑这样的结构类，它没有辅助基集，只有一个主基集A，其典型特征为 \(V \in \mathfrak{P}(\mathfrak{P}(A))\) ，以及作为公理的可迁移关系
\end{enumerate}

\[
\begin{array}{l} (\forall \mathrm{V}) ((\mathrm{V} ^ {\prime} \subset \mathrm{V}) \Longrightarrow ((\bigcup_ {\mathrm{X} \in \mathrm{V} ^ {\prime}} \mathrm{X}) \in \mathrm{V})) \\ \text { 且 } \quad (\forall \mathrm{X}) (\forall \mathrm{Y}) ((\mathrm{X} \in \mathrm{V} \text {   且   } \mathrm{Y} \in \mathrm{V}) \Longrightarrow ((\mathrm{X} \cap \mathrm{Y}) \in \mathrm{V})). \end{array}
\]

这种结构类型被称为拓扑结构类型。这种类型的一个结构也称为拓扑，而该关系 \(X \in V\) 用拓扑V中的开集X来表示（一般拓扑学，第一章，§1）。

* (4) 取 \(\mathfrak{G}\) 作为除环结构这个物种的理论，其（除其他事项外）有一个常数K作为唯一的（主）基集。K上的左向量空间这一结构物种以K为辅助基集，以E为主基集，其典型特征关系为

\[
\mathbf {V} \in \mathfrak {P} ((\mathbf {E} \times \mathbf {E}) \times \mathbf {E}) \times \mathfrak {P} ((\mathbf {K} \times \mathbf {E}) \times \mathbf {E})
\]

\((pr_{1}V \text{ 是 } E \text{ 中加法的图，而 } pr_{2}V \text{ 是标量乘法的图 }); \text{ 我们在此不陈述该结构物种的公理。}\)

\begin{enumerate}
\def\labelenumi{(\arabic{enumi})}
\setcounter{enumi}{4}
\tightlist
\item
  再次令 \(\mathfrak{G}\) 是集合论；在这个理论中，复数域 \(\mathbf{C}\) 是一个不含字母的项。维数为 \(n\) 的复解析流形的结构种具有 \(\mathbf{C}\) 作为辅助基集，以及一个主基集 \(V\) 我们在此不指出这类结构的典型特征或公理。*
\end{enumerate}

\minorheading{注记}

\begin{enumerate}
\def\labelenumi{(\arabic{enumi})}
\tightlist
\item
  在应用中，常常出现（如上述例4中）梯队 \(S(E_1, \ldots, E_n, A_1, \ldots, A_m)\) 是若干梯队的乘积。
\end{enumerate}

\[
\mathrm{S} _ {1} \left(\mathrm{E} _ {1}, \dots , \mathrm{A} _ {m}\right) \times \dots \times \mathrm{S} _ {p} \left(\mathrm{E} _ {1}, \dots , \mathrm{A} _ {m}\right).
\]

若是如此，在 \(s\) 的定义中，字母 \(\Sigma\) 常被替换为一个“ \(p\) 元组” \((s_{1}, \ldots, s_{p})\) （参见第二章，第2节，第1号）。

此外，结构种类的公理 \(\Sigma\) 通常被写成若干可迁移关系的合取（如上面的例3所示）。这些关系被称为该物种的公理。 \(\Sigma\) .

\begin{enumerate}
\def\labelenumi{(\arabic{enumi})}
\setcounter{enumi}{1}
\item
  数学中最常用的结构种类，以及赋予这些结构种类的集合，都有特定的名称。因此，有序集（第三章，§1）是赋予序结构（例1）的集合；* 在本系列后续的卷册中，我们将定义群、域、拓扑空间、微分流形等概念，所有这些都指代赋予特定结构的集合。*
\item
  按照惯例，在集合论中 \(\mathfrak{G}\) ，赋予 \(n\) 不同的字母 \(x_{1},\ldots,x_{n}\) （没有典型刻画且没有公理）被视为结构的一种 \(\Sigma_{0}\) ，称为集合上的结构， \(n\) 主基集 \(x_{1},\ldots,x_{n}\) .
\end{enumerate}

